\chapter{Greater Trochanteric Pain Syndrome (GTPS)}

\section*{Definition and Overview}
Greater trochanteric pain syndrome (GTPS) is a common cause of pain on the outside of the hip and thigh. It refers to a group of conditions affecting the tissues around the greater trochanter, the bony prominence on the outside of the upper thigh bone, most commonly the gluteal tendons and the bursa that cushions them. GTPS was formerly called trochanteric bursitis, but it is now recognised that the main problem is usually tendinopathy of the gluteal muscles rather than bursitis alone. It causes aching or sharp pain on the outer hip that worsens with activity, lying on the affected side, or climbing stairs. Most people improve with physiotherapy, activity modification, and pain management.

\section*{Epidemiology}
GTPS is a very common cause of hip pain, affecting around 15--25\% of adults at some stage, and it becomes more common with age. It is most frequent in women, particularly those aged 40--70, and is associated with obesity, leg length differences, hip osteoarthritis, and walking patterns that strain the gluteal tendons. It is common in runners and in people who do a lot of walking, stair climbing, or standing. Although it is often chronic and can be disabling, most people respond well to conservative treatment.

\section*{Aetiology and Pathophysiology}
GTPS results from overload and degeneration of the gluteal tendons and the surrounding bursa.

\begin{itemize}
    \item \textbf{Anatomy:} the gluteus medius and minimus tendons attach to the greater trochanter and are the main stabilisers of the hip during walking; the trochanteric bursa lies between them and the bone
    \item \textbf{Tendinopathy:} repetitive load, weakness, or altered biomechanics cause degenerative changes in the tendons, with pain and reduced capacity
    \item \textbf{Bursitis:} the bursa can become inflamed in response to tendon friction, though this is often a secondary feature
    \item \textbf{Biomechanical factors:} weakness of the gluteal muscles, poor control of the pelvis, leg length difference, and foot problems
    \item \textbf{Contributing conditions:} hip osteoarthritis, lumbar spine problems, and prior hip surgery
    \item \textbf{Load factors:} sudden increases in walking, running, or stair climbing, and prolonged standing on hard surfaces
\end{itemize}

\section*{Clinical Features}
The pain of GTPS is characteristic in its location and behaviour.

\begin{itemize}
    \item Pain on the outside of the hip that may spread down the outer thigh
    \item Pain when lying on the affected side, which disturbs sleep
    \item Pain with walking, running, climbing stairs, or standing for long periods
    \item Tenderness over the greater trochanter when pressed
    \item Pain when standing on one leg or rising from a chair
    \item Stiffness, especially after sitting for a long time
    \item A feeling of weakness or giving way in the hip
    \item Symptoms that are worse with activity and better with rest, though they can be persistent
\end{itemize}

\section*{Differential Diagnosis}
Hip pain has many causes, and accurate diagnosis is important.

\begin{itemize}
    \item Hip osteoarthritis, which causes pain in the groin and front of the hip, and stiffness
    \item Lumbar spine conditions with referred pain into the buttock and thigh
    \item Sacroiliac joint dysfunction
    \item Femoroacetabular impingement and labral tears in younger people
    \item Femoral neck stress fracture and other fractures
    \item Avascular necrosis of the femoral head
    \item Sciatica and other nerve-related pain
    \item Fibromyalgia and other widespread pain conditions
\end{itemize}

\section*{When to Seek Medical Care}
See a doctor or physiotherapist if hip pain persists for more than a few weeks, is severe, or limits daily activities and sleep. Seek early assessment if the pain follows a fall or injury.

\textbf{Call 000 (or local emergency services) for:}
\begin{itemize}
    \item Sudden severe hip or thigh pain with inability to bear weight, which may indicate a fracture
    \item Hip pain with fever, which may indicate a joint infection
    \item Sudden leg swelling, redness, or chest symptoms, which may indicate a blood clot (deep vein thrombosis)
    \item Inability to move the leg or foot
\end{itemize}

\section*{Diagnosis and Investigations}
GTPS is usually diagnosed clinically, with imaging used when the diagnosis is uncertain or treatment fails.

\begin{itemize}
    \item \textbf{Clinical examination:} tenderness over the greater trochanter, reproduced by specific tests such as the resisted hip abduction test
    \item \textbf{History:} the location and behaviour of the pain are highly suggestive
    \item \textbf{Ultrasound:} can show tendon thickening, tears, and bursal fluid
    \item \textbf{MRI:} used for persistent cases or to exclude other causes, showing tendon and bone detail
    \item \textbf{X-rays:} rule out arthritis and fractures
    \item \textbf{Response to treatment:} improvement with physiotherapy supports the diagnosis
\end{itemize}

\section*{Management}
Conservative treatment is effective in most cases and is the first-line approach.

\begin{itemize}
    \item \textbf{Physiotherapy:} the cornerstone of treatment, focusing on strengthening the gluteal muscles and improving hip and pelvic control
    \item \textbf{Activity modification:} reducing aggravating activities such as long walks, stairs, and running, then gradually returning
    \item \textbf{Pain relief:} paracetamol or anti-inflammatory medicines for short-term symptom control
    \item \textbf{Ice:} applying ice to the outer hip after activity
    \item \textbf{Footwear and orthotics:} addressing leg length differences and foot issues
    \item \textbf{Injections:} corticosteroid injections can provide short-term relief and allow physiotherapy to proceed
    \item \textbf{Shockwave therapy:} extracorporeal shockwave therapy is used for persistent tendinopathy
    \item \textbf{Surgery:} rarely needed, but tendon repair may be considered when conservative treatment fails
\end{itemize}

\section*{Complications}
\begin{itemize}
    \item Chronic pain and disability affecting mobility, work, and sleep
    \item Altered walking pattern that strains the lower back, knee, and opposite hip
    \item Tendon tears, which occur in some cases of gluteal tendinopathy
    \item Side effects of injections, including weakening of the tendon with repeated steroid use
    \item Reduced physical activity leading to deconditioning and weight gain
    \item Complications of surgery in the small minority who need it
\end{itemize}

\section*{Prognosis and Outcomes}
Most people with GTPS improve with conservative treatment. Physiotherapy is effective in around 70--80\% of people, though improvement can take several months, and symptoms may flare with increased activity. Corticosteroid injections provide faster short-term relief but are not a long-term cure. About 10--20\% of people have persistent symptoms despite treatment, and a small proportion need surgery. The best outcomes come from a structured exercise program combined with gradual return to activity.

\section*{Prevention and Self-Care}
\begin{itemize}
    \item Strengthen the gluteal and hip muscles with regular, progressive exercise
    \item Maintain a healthy weight to reduce load on the hip
    \item Increase walking, running, and exercise loads gradually
    \item Avoid prolonged standing on hard surfaces, and take breaks
    \item Improve your walking and running technique with professional advice
    \item Sleep with a pillow between the knees to reduce pressure on the hip
    \item Treat underlying problems such as leg length differences with appropriate footwear
    \item Start physiotherapy early rather than waiting for symptoms to settle
    \item Resume activity gradually after treatment to avoid relapse
\end{itemize}

\section*{Sources and Further Reading}
The following reputable sources provide further information for healthcare students and patients seeking reliable, current advice about greater trochanteric pain syndrome.

\begin{itemize}
    \item Sports Medicine Australia. \textit{Hip pain fact sheet.} https://sma.org.au/
    \item Australian Physiotherapy Association. \textit{GTPS and hip pain information.} https://australian.physio/
    \item Healthdirect Australia. \textit{Hip pain and greater trochanteric pain syndrome.} https://www.healthdirect.gov.au/
    \item Arthritis Australia. \textit{Hip pain and its causes.} https://arthritisaustralia.com.au/
    \item NHS. \textit{Hip pain: causes and treatment.} https://www.nhs.uk/
    \item American Academy of Orthopaedic Surgeons. \textit{Trochanteric bursitis and gluteal tendinopathy.} https://orthoinfo.aaos.org/
\end{itemize}
